\section{中文编译（一）：旅程开场与总体感受}

\begin{translationbox}{原文主线编译}
作者从杭州到上海的高铁上写下开场：窗外是山脊、风机、田野和高楼混杂的景象。他说自己带着很大的谦逊离开中国。到一个陌生地方却被热情接待，是一种非常温暖、非常有人味的体验。他见到了许多过去只在远处知道的 AI 生态参与者，而这些人用笑容和善意欢迎他，让他重新意识到：自己的工作和 AI 生态本身都是全球性的。
\end{translationbox}

这个开场很重要。作者不是从 benchmark、芯片制裁或地缘政治开始，而是从人的温度开始。后文虽然讨论技术和产业竞争，但语气并不冷冰冰。他把中国研究者写成具体的人：谦逊、专注、热情、有执行力。这种写法本身就在反抗一种常见叙事：把中国 AI 实验室只当作美国政策、芯片限制或产业竞争中的抽象对象。

\begin{knowledgebox}{为什么开场的``人味''重要}
AI 产业讨论很容易被压缩成国家、公司、模型榜单和政策。作者的开场提醒读者：前沿 AI 竞争背后是具体研究者、学生、工程师和团队的日常工作。理解这些人的动机和组织环境，往往比只看参数量和 benchmark 更接近真实。
\end{knowledgebox}

\subsection{本章小结}

文章一开始就奠定了语气：这不是一篇冷战式中国 AI 分析，而是一篇带着尊重、好奇和不确定性的现场笔记。


